%%=============================================================================
%% Methodologie
%%=============================================================================

\chapter{\IfLanguageName{dutch}{Methodologie}{Methodology}}%
\label{ch:methodologie}

%% TODO: In dit hoofstuk geef je een korte toelichting over hoe je te werk bent
%% gegaan. Verdeel je onderzoek in grote fasen, en licht in elke fase toe wat
%% de doelstelling was, welke deliverables daar uit gekomen zijn, en welke
%% onderzoeksmethoden je daarbij toegepast hebt. Verantwoord waarom je
%% op deze manier te werk gegaan bent.
%% 
%% Voorbeelden van zulke fasen zijn: literatuurstudie, opstellen van een
%% requirements-analyse, opstellen long-list (bij vergelijkende studie),
%% selectie van geschikte tools (bij vergelijkende studie, "short-list"),
%% opzetten testopstelling/PoC, uitvoeren testen en verzamelen
%% van resultaten, analyse van resultaten, ...
%%
%% !!!!! LET OP !!!!!
%%
%% Het is uitdrukkelijk NIET de bedoeling dat je het grootste deel van de corpus
%% van je bachelorproef in dit hoofstuk verwerkt! Dit hoofdstuk is eerder een
%% kort overzicht van je plan van aanpak.
%%
%% Maak voor elke fase (behalve het literatuuronderzoek) een NIEUW HOOFDSTUK aan
%% en geef het een gepaste titel.

Dit onderzoek volgt een Proof-of-Concept-aanpak, een veelgebruikte onderzoeksbenadering binnen toegepaste informatica waarbij een concreet prototype wordt ontwikkeld en geëvalueerd in functie van een reële praktijksituatie.
De centrale vraag is in welke mate een interactieve interface huisartsen en praktijkmedewerkers kan helpen om toegankelijkheidsevaluaties sneller, consistenter en gebruiksvriendelijker uit te voeren.
Om deze vraag te beantwoorden wordt het onderzoek opgedeeld in vier opeenvolgende fasen: een ontwerpfase, de ontwikkeling van de interface, een empirische evaluatie en ten slotte de analyse en rapportering van de resultaten.
\subsection{Fase 1: Ontwerpfase en validatie via de werkgroep}
\label{ter:Ontwerpfase}

De eerste fase bestaat uit de ontwerpfase, waarin op basis van de bestaande indicatorensets en de literatuur een eerste visueel prototype van de interface werd uitgetekend in Figma. 
De keuze voor Figma als ontwerptool is gemotiveerd door de mogelijkheid om snel een klikbaar, interactief prototype te realiseren zonder enige implementatie, 
waardoor het ontwerp vroeg in het proces aan de werkgroep kon worden voorgelegd. 

Het Figma-prototype werd vervolgens gepresenteerd aan de werkgroep van Eerstelijnszone Aalst, 
bestaande uit huisartsen, onderzoekers en medewerkers van Huisartsenwachtpost regio Aalst die nauw betrokken zijn bij de toegankelijkheidsproblematiek. 
Tijdens deze presentatie werd samen met de werkgroep het prototype doorlopen. 
En werd feedback gegeven op het ontwerp en de gebruiksvriendelijkheid. Deze aanpak laat toe om vroegtijdig te valideren of het voorgestelde concept aansluit bij de noden van de doelgroep, zonder dat hiervoor reeds een werkende implementatie vereist is. 
De verzamelde feedback werd kwalitatief geanalyseerd en vertaald naar een set concrete aanpassingen aan het ontwerp, die als input dienden voor de ontwikkelfase.
De output van deze fase bestaat uit een bijgewerkt Figma-prototype dat de verwerkte feedback weerspiegelt en als blauwdruk dient voor de technische implementatie in de volgende fase.

\subsection{Fase 2: Implementatie}
\label{ter:Implementatie}
In de tweede fase wordt het goedgekeurde Figma-ontwerp technisch geïmplementeerd. 
De frontend wordt gerealiseerd in Angular, met SurveyJS als bibliotheek voor de vragenlijstlogica. 
Voor de backend wordt gebruikgemaakt van Express als Node.js-framework, Prisma als ORM en PostgreSQL als relationele databank. 
Deze technologiekeuzes zijn gemotiveerd door hun onderlinge compatibiliteit, de brede community-ondersteuning en de mogelijkheid om de volledige stack lokaal te draaien, wat aansluit bij de privacyvereisten die inherent zijn aan de verwerking van gegevens over huisartsenpraktijken.

De applicatie wordt na implementatie gehost via GitHub Pages, wat een laagdrempelige en gratis hostingoplossing biedt voor statische frontend-applicaties. 
De backend wordt apart gedeployed zodat de API-aanroepen vanuit de Angular-applicatie correct kunnen worden afgehandeld. 
Deze fase resulteert in een volledig werkende en publiek toegankelijke proof-of-concept interface.
\subsection{Fase 3: Gebruikerstesting en evaluatie}
\label{ter:testing}
De derde fase omvat de gebruikerstesting, waarin de gehoste applicatie wordt voorgelegd aan een groep huisartsen en praktijkmedewerkers. 
De deelnemers doorlopen de interface zelfstandig en vullen na afloop een vragenlijst in die twee componenten bevat: de System Usability Scale (SUS) voor een gestandaardiseerde meting van de ervaren gebruiksvriendelijkheid, en een open feedbackgedeelte waarin deelnemers vrij kunnen aangeven wat goed werkte en wat voor verbetering vatbaar is.
De SUS-score wordt berekend op basis van de tien gestandaardiseerde stellingen en gepositioneerd ten opzichte van de internationaal erkende drempelwaarde van 75, die geldt als indicatie voor goede usability. 
De vrije feedback wordt kwalitatief geanalyseerd en gegroepeerd per thema, zodat terugkerende knelpunten of positieve bevindingen duidelijk naar voren komen. 
De combinatie van beide meetmethoden laat toe om zowel een kwantitatief oordeel als een inhoudelijk genuanceerd beeld te vormen van de gebruikservaring.
\subsection{Fase 4: Analyse en rapportering}
\label{ter:analyse}
In de vierde fase worden de verzamelde gegevens geanalyseerd en geïnterpreteerd. 
De SUS-scores worden berekend en gepositioneerd ten opzichte van de drempelwaarde van 75. 
De kwalitatieve feedback van de deelnemers wordt thematisch gecodeerd om terugkerende patronen te identificeren. 
De combinatie van beide datatypen maakt het mogelijk om een genuanceerd antwoord te formuleren op de centrale onderzoeksvraag en om te beoordelen welke onderdelen van de interface effectief bijdragen aan efficiëntie en gebruiksvriendelijkheid, 
en welke nog verdere verfijning vereisen.